\subsection{CLOSED SETS}

DEFINITION 7. In a topological space X, the complements of the open sets of X are called closed sets.

Taking complements, we find that the axioms \((\mathbf{O}_{\mathbf{I}})\) and \((\mathbf{O}_{\mathbf{II}})\) take the following form:

\((\mathbf{O}_{\mathbf{I}}^{\prime})\) Every intersection of closed sets is a closed set.

\((\mathbf{O}_{\mathrm{II}}^{\prime})\) Every finite union of closed sets is a closed set.

The empty set and the whole space X are closed (and therefore both open and closed; cf.~§ 11).

On the rational line, every interval of the form \([a, \rightarrow[\) is a closed set, for its complement \(]\leftarrow, a[\) is open; likewise, every interval of the form \(]\leftarrow, a]\) is a closed set; hence so is every bounded closed interval \([a, b]\) , since it is the intersection of the intervals \([a, \rightarrow[\) and \(]\leftarrow, b]\) .

The set Z of rational integers is closed in the rational line, since its

complement \(\bigcup_{n\in \mathbf{Z}}]n,n + 1[\) is open.

A covering \((\mathbf{F}_i)_{i\in I}\) of a subset \(\mathbf{A}\) of a topological space \(\mathbf{X}\) is said to be closed if each of the \(\mathbf{F}_i\) is closed in \(\mathbf{X}\) .

A homeomorphism \(f\) of a topological space \(\mathbf{X}\) onto a topological space \(\mathbf{X}'\) (no. 1) can be characterized as a bijection of \(\mathbf{X}\) onto \(\mathbf{X}'\) such that the image under \(f\) of every closed subset of \(\mathbf{X}\) is a closed subset of \(\mathbf{X}'\) and the inverse image under \(f\) of every closed subset of \(\mathbf{X}'\) is a closed subset of \(\mathbf{X}\) .

